\exnumber{00}
\extitle{UX VI - Ex 1 - User flow}
\exfiles{ex00/Figma file, named Name\_FirstName\_DeliverableName\_Date\_VersionNumber}
\exauthorize{Those named in the exercise}
\makeheaderfiles

\textbf{Context:}

Your company's client is a movie theaters chain called, Movie 4 all.

They want to launch a platform for people who want to watch movies but don't want to spend endless time selecting one.
You are replacing someone in the company who has already worked on the user research, the problem statement and the list of features they plan to develop. Your goal is to prototype and test these features.

\textbf{Instructions:}

Based on the 3 features that have come out of the ideation workshops led by your collegues, design 3 user flows (1 per feature).

Features:

\begin{itemize}
\item 1: Adding a filter based on their current mood
\item 2: Offer a selection of 3 movies to be watched
\item 3: Ask users to rate the movie after watching it for further recommendation
\end{itemize}

\textbf{Deliverables:}

\begin{itemize}
\item 3 user flows on a board (digital or pictures of a paper board)
\item Upload the documents on Github, in a zip folder named "ProjectTitle\_Name\_FirstName", with all project deliverables as follows: "Name\_FirstName\_DeliverableName\_Date\_VersionNumber".  For example, the first deliverable can be named as "Doe\_John\_MidFiPrototype\_05242024\_V1".
\item Add a title within all written documents.
\item Pay attention to the global aspect of the deliverables. They must be clear, simple and easy to read. You can get inspiration on canva.com but don't overload your design with too much details!
\end{itemize}

\textbf{Resources:}

\begin{itemize}
\item User flow glossary \url{https://www.productplan.com/glossary/user-flow/}
\item Site flows vs User Flows \url{https://uxmovement.com/wireframes/site-flows-vs-user-flows-when-to-use-which/}
\item How to make a User Flow diagram \url{https://www.lucidchart.com/blog/how-to-make-a-user-flow-diagram}
\end{itemize}

\textbf{Recommended tools:}

\begin{itemize}
\item Figma
\item Whimsical
\end{itemize}

